# Profil J2 — Analyste, 6 mois d'ancienneté

**Fonction :** Analyste  
**Ancienneté :** 6 mois  
**Missions :** Gestion de crise en entreprise, réponses à appels d'offre, préparation d'exercices (PowerPoint, Word), restitution et analyse post-exercice (Word, Excel, PowerPoint)  
**Rapport à l'IA :** Indifférent — utilise pour aller plus vite et trouver des idées d'injects  
**Tic de langage :** "En vrai"  
**Contexte spécifique :** Vient d'un stage dans l'armée où l'IA était interdite et mal vue. S'acclimate bien et utilise Copilot régulièrement par pragmatisme.  
**Outils utilisés :** Copilot officiel uniquement (pas d'officieux)

---

## T0 — Cadrage (parcours et missions)

> "Je suis analyste ici depuis six mois, après un stage dans l'armée où j'ai bossé sur des exercices de gestion de crise. En vrai, c'est ce qui m'a amené ici — je me suis rendu compte que j'aimais bien les situations où il faut réagir vite et structurer l'incertitude. Ici, je fais beaucoup de réponses à appels d'offre, donc des Powerpoints de propale. Ensuite, sur les exercices de crise eux-mêmes, je prépare des présentations pour les équipes, je fais des documents Word pour la préparation, puis le jour J, un PowerPoint avec le déroulé exact. Après l'exercice, on fait des analyses — donc Word, Excel, et souvent une restitution en PowerPoint aussi. C'est un cycle assez complet."

---

## T1 — Usage réel (Appropriation)

### Ouverture — dernière fois que vous avez utilisé un outil génératif

> "La dernière fois, c'était pour une propale. On avait trois jours pour répondre, et je devais structurer une section sur les 'leviers d'innovation' dans un exercice de crise. J'ai ouvert Copilot dans PowerPoint, je lui ai filé le contexte en deux phrases, et je lui ai demandé de me sortir cinq pistes d'injects innovants — des événements imprévus qu'on pourrait intégrer pour rendre l'exercice plus réaliste. Il m'a sorti des trucs genre 'panne simultanée des systèmes de communication' ou 'fuite d'information dans la presse'. Certaines idées étaient utilisables, d'autres trop génériques. J'ai gardé trois pistes, j'ai adapté les formulations, et j'ai construit la slide derrière. En vrai, ça m'a fait gagner du temps sur la phase de brainstorming."

### Relance — sur quelles tâches au quotidien ?

> "Je l'utilise surtout pour deux choses : les propales et les idées d'injects. Pour les propales, c'est comme je viens de dire — il m'aide à structurer et à trouver des angles. Pour les injects, c'est plus créatif : je lui demande des scénarios, des rebondissements, des données fictives mais crédibles. Par contre, je ne l'utilise jamais sur les documents d'analyse post-exercice. Là, c'est du chiffrage, des retours d'expérience concrets, des recommandations ancrées dans ce qui s'est passé. L'IA n'a pas été dans la salle, elle ne sait pas ce qui s'est vraiment joué. Donc je fais tout à la main."

### Relance — comment as-tu appris à t'en servir ?

> "J'ai appris sur le tas. Le cabinet a mis Copilot à disposition, mais il n'y a pas eu de formation. J'ai juste commencé à l'utiliser sur des petites tâches pour voir ce que ça donnait. En vrai, c'est assez intuitif — tu poses une question, il te répond. Je ne suis pas allé chercher des tutoriels ou des prompts avancés. Je fais simple."

### Relance — le temps gagné, tu l'as mis sur quoi ?

> "Sur la propale que je viens de mentionner, j'ai gagné une bonne heure sur la phase de recherche d'idées. Je l'ai utilisée pour affiner ma structure et pour vérifier que mes arguments tenaient la route. Donc j'ai passé plus de temps sur la cohérence globale du dossier plutôt que sur le remplissage des slides. C'est un bon équilibre."

### Relance — une situation où ça t'a demandé plus d'effort que de le faire toi-même ?

> "Une fois, j'ai voulu lui faire générer un inject hyper spécifique sur une crise cyber dans le secteur financier. Il m'a sorti des trucs tellement génériques que c'était inutilisable. Des 'attaques DDoS' et des 'ransomwares' sans aucune originalité. J'ai passé dix minutes à reformuler mon prompt, à lui donner des exemples, à essayer de le guider. Rien n'y faisait. Au final, j'ai écrit l'inject moi-même en cinq minutes. Depuis, je sais que pour des sujets trop spécifiques, il vaut mieux utiliser sa tête."

### Relance — vous en parlez entre vous ?

> "Oui, un peu. Avec les autres analystes, on se dit parfois 't'as essayé Copilot pour ça ?' ou 'il t'a sorti quoi ?'. Mais c'est assez informel. C'est pas un sujet qu'on aborde en réunion d'équipe ou avec les managers. En vrai, j'ai l'impression que tout le monde l'utilise un peu, mais personne n'en parle ouvertement. C'est un peu comme si c'était un secret de polichinelle."

---

## T2 — Apprentissage du métier (Compétences)

### Ouverture — comment on apprend le métier aujourd'hui ?

> "Je pense que j'apprends surtout par la pratique. Les propales, les exercices, les restitutions — c'est en faisant qu'on comprend ce qui fonctionne ou pas. L'IA m'aide à aller plus vite, mais elle ne m'apprend pas le métier. Par exemple, sur la gestion de crise, c'est en voyant comment les clients réagissent, en écoutant les retours des seniors, en analysant ce qui a bien ou mal marché qu'on apprend. L'IA ne peut pas remplacer cette expérience de terrain. En vrai, c'est un outil, pas un mentor."

### Relance — par quelles tâches on apprend réellement le métier ?

> "Les tâches où on est confronté à l'incertitude. Une propale, tu ne sais jamais exactement ce que le client va attendre. Un exercice de crise, tu as des paramètres qui changent en cours de route. C'est là que tu progresses — dans l'adaptation. Les tâches répétitives, comme la mise en forme ou la reformulation, elles n'apprennent rien. L'IA les fait très bien, et tant mieux. Mais le cœur du métier, c'est la prise de décision sous pression, et ça, l'IA ne le fait pas."

### Relance — les tâches ingrates, comme les injects, sont-elles encore formatrices ?

> "Les injects, c'est un cas intéressant. Les trouver soi-même, c'est un exercice créatif. Mais je trouve que l'IA peut être un bon partenaire de brainstorming. Elle te sort des pistes, tu les évalues, tu les améliores. Ça te fait gagner du temps sur la génération, et tu consacres ce temps à l'affinage et à la crédibilité. Donc je ne dirais pas que l'IA appauvrit cet apprentissage — elle le rend plus efficace. Mais il faut garder la main sur le choix final."

### Relance — comment vous vérifiez ce que produit l'outil ?

> "Je vérifie surtout la cohérence globale. Est-ce que ce qu'il me propose est crédible dans le contexte de l'exercice ? Est-ce que les idées s'emboîtent bien avec ce qu'on sait du client ? Si je vois des incohérences ou des trucs trop génériques, je les écarte ou je les retravaille. En vrai, je ne vais pas vérifier chaque source ou chaque chiffre, parce que je ne l'utilise pas sur des données factuelles. C'est plus pour la créativité et la structure."

---

## T3 — Client et valeur (Légitimité de l'expert)

### Ouverture — la confiance du client repose sur quoi ?

> "La confiance du client, à mon avis, repose sur notre capacité à anticiper. Dans un exercice de crise, le client veut voir qu'on a pensé à tout — les scénarios, les imprévus, les conséquences. Si on arrive avec un plan solide et qu'on est capable de s'adapter en direct, la confiance est là. L'IA peut aider à préparer, mais le jour J, c'est le consultant qui est sur le terrain. C'est lui qui doit être crédible."

### Relance — diriez-vous à un client qu'une partie a été produite par l'IA ?

> "Non, je ne pense pas. Pas parce que je cache quelque chose, mais parce que ça n'apporte rien au client. Ce qu'il veut, c'est un livrable de qualité. Peu importe comment on y est arrivé. Si je lui dis 'j'ai utilisé Copilot', je risque de créer une distraction. Il va se demander si c'est vraiment nous qui avons fait le boulot. Donc je ne mentionne pas l'outil."

### Relance — si un livrable était perçu comme moins crédible, pourquoi ?

> "Ce serait probablement une question d'ancrage. Si le client sent qu'on n'a pas compris son secteur, ses contraintes, sa culture, il va douter. L'IA peut générer du contenu générique, mais elle ne peut pas saisir les nuances du contexte. Donc le risque, c'est un livrable qui sonne juste mais qui n'est pas vraiment pertinent. C'est pour ça que je garde la main sur tout ce qui touche à l'analyse et au jugement."

### Relance — dans dix ans, sur quoi un client acceptera-t-il de payer un cabinet ?

> "Je pense qu'il paiera pour l'expérience humaine. La gestion de crise, c'est du relationnel, de l'improvisation, de la prise de risque. Une IA peut préparer, elle peut même simuler. Mais elle ne peut pas être dans la salle, sentir la tension, lire les visages, ajuster son discours en direct. Le conseil de demain, ce sera une affaire de présence, pas d'algorithmes."

---

## T4 — Gouvernance et risques

### Ouverture — il existe des règles sur l'usage des outils ici ?

> "Oui, on sait que Copilot est l'outil officiel. On a eu une petite communication sur le sujet, mais c'était assez général — genre 'utilisez-le à bon escient'. Rien de très contraignant. Après, avec mon passé dans l'armée, j'ai une sensibilité particulière à la confidentialité. Je n'ai pas besoin qu'on me dise de faire attention aux données sensibles, je le fais naturellement. Donc les règles officielles, je les connais, mais elles ne changent pas grand-chose à ma pratique."

### Relance — avant qu'un livrable ne parte chez le client, il se passe quoi ?

> "Il y a une relecture par le manager ou le senior. Ils vérifient la cohérence, les chiffres, les recommandations. Mais ils ne demandent pas spécifiquement si on a utilisé l'IA. Je pense qu'ils s'en doutent parfois, mais ils ne posent pas la question. Tant que le contenu est bon, ils ne regardent pas le processus."

### Relance — déjà rencontré un contenu produit par l'IA qui s'est révélé faux ?

> "Pas vraiment. Des trucs génériques, oui, mais pas de l'inventé pur et dur. Après, je l'utilise surtout pour des idées, pas pour des faits. Donc le risque d'hallucination est limité dans mon usage. Mais j'ai déjà vu des collègues utiliser ChatGPT pour des données et se retrouver avec des chiffres qui sortaient de nulle part. Ils ont dû tout reprendre. En vrai, je préfère ne pas prendre ce risque."

### Relance — si vous écriviez la règle vous-même, vous diriez quoi ?

> "Je dirais : 'L'IA est un outil d'aide à la créativité et à la structuration, pas un substitut au jugement.' Je mettrais l'accent sur la vérification — si tu utilises l'IA pour des faits, tu dois être capable de les justifier. Et j'ajouterais une règle sur les données sensibles : pas de données client dans les outils non-officiels. Mais en vrai, ce n'est pas compliqué à appliquer."

---

## T5 — Clôture

### Relance — à la place de la direction, quelle serait votre première décision ?

> "Je formaliserais une charte simple. Pas un document de vingt pages, mais une page avec trois règles : ce qu'on peut faire, ce qu'on ne peut pas faire, et ce qu'on doit absolument vérifier. Et je mettrais en place un partage de bonnes pratiques entre juniors — un canal Teams ou une réunion informelle. Parce qu'aujourd'hui, on apprend en tâtonnant, et certains font des erreurs qu'ils auraient pu éviter."

### Relance — un aspect important qu'on n'a pas abordé ?

> "La culture du cabinet. Je viens de l'armée, où l'IA n'était pas vue d'un bon œil. Ici, c'est plus ouvert, mais il y a quand même une ambivalence. On nous dit d'innover, mais on ne nous dit pas clairement comment intégrer ces outils. En vrai, il y a un décalage entre le discours et la pratique."

### Relance — quelqu'un dont le point de vue serait particulièrement utile ?

> "Un consultant qui a dix ans d'expérience dans la gestion de crise, pour voir comment il perçoit ces outils. Est-ce qu'il les utilise ? Est-ce qu'il les craint ? Est-ce qu'il voit une différence entre les juniors qui les utilisent et ceux qui n'en font pas ? Je pense que son regard serait éclairant."